\section{Split-pair anchor: cyclic cubic}\label{sec:cyclic}

We treat the last residual image by an Eisenstein comparator.  The
integral calculation has two parts: a primitive ray-unit determinant
and its realization as a pair of CM-disk logarithms.  Both use full
character-weighted sums, so the comparison retains their dyadic
indices.

\subsection{The Selmer seed and the comparator}

Let \(E_0/\mathbb Q\) be non-CM with residual image \(C_3\), and put
\(W=E_0[2]\).  Over \(\mathbb F_4\), the representation splits into
characters lifted by \(\rho,\rho^{-1}\), where \(\rho\) is a primitive
even Dirichlet character of order three and odd conductor \(m\).
We use arithmetic character conventions and coefficients
\(\mathcal O=\mathbb Z_2[\mu_3]\).

\begin{lemma}[Cyclic-cubic Selmer seed]\label{cyclic:seed}
There are a prime \(d\equiv7\pmod8\), a field
\(K=\mathbb Q(\sqrt{-d})\), and a positive fundamental discriminant
\(h\), allowing \(h=1\), such that:
\begin{enumerate}[label=\textup{(\roman*)}]
\item every prime dividing \(2mN_{E_0}\) splits in \(K\), and
      \(\rho(d)\ne1\);
\item \(h\) is a product of fresh primes congruent to one modulo eight,
      good for \(E_0\), split in \(K\), and prime to \(d\);
\item the finite two-Selmer dimensions of \(E_0^h\) and
      \(E_0^{-dh}\) are zero and one, in some order.  In particular,
      \(s_2(E_0^h/K)=1\).
\end{enumerate}
The field \(K\) has only the units \(\pm1\) and has odd class number.
\end{lemma}

\begin{proof}
The cubic prescription \(\rho(d)\ne1\) is independent of the
quadratic splitting prescriptions, so Chebotarev supplies \(d\).
The discriminant of \(K\) is \(-d\), with a single prime factor;
genus theory gives odd class number.  The congruence at eight gives
splitting at two.

The two twists \(E_0,E_0^{-d}\) have opposite Selmer parities, as in
the imaginary \(S_3\) construction.  Twist simultaneously at fresh
good primes which are local squares at the old places, split in
\(K\), and split on \(W\).  The ramified Kummer condition there is
transverse to the old finite condition: it is obtained by taking halves
of invariant two-torsion.  Reciprocity shows that a switch lowers the
Selmer dimension by at least two when the old finite evaluation has
rank two, and does not increase it when that evaluation has rank one.

We can make these prescriptions simultaneously.  The space \(W\) is a
line over its endomorphism field \(\mathbb F_4\).  On the joint kernel
of \(W\) and the abelian congruence data, evaluations test arbitrary
linear maps on the \(\mathbb F_4\)-span of the chosen cohomology
classes.  Restriction detects the classes because of the nontrivial
central scalar.  The joint evaluations span, and additivity together
with cubic equivariance makes the actual image an
\(\mathbb F_4\)-subspace.

For two \(\mathbb F_2\)-independent classes, a rank-two evaluation is
possible even if they are proportional over \(\mathbb F_4\): their
ratio then lies outside \(\mathbb F_2\).  If they are independent over
\(\mathbb F_4\), either of the two ratios in
\(\mathbb F_4\setminus\mathbb F_2\) gives rank two.  At most one of
these ratios kills a specified third nonzero class in their span; if
that class is outside the span, prescribe its evaluation separately.
We may therefore cut the odd Selmer dimension to one and then the even
dimension to zero, keeping nonzero evaluation on the odd line during
the second stage.  Chebotarev realizes the required evaluations with
all the stated local square conditions.  Parity, the absence of
rational two-torsion, and quadratic induction now give
\(s_2(E_0^h/K)=1\).
\end{proof}

For the rest of this section put \(E=E_0^h\).  In particular, \(K\)
splits every prime dividing \(2N_E\).  No analytic-rank assertion is
being made at this stage.

Retain the true free \(t\)-direction and several order-conductor variables
\(u_j\) at good inert primes \(r_j\to-1\) two-adically, independent cyclic
Sylow quotients as before. Write \(\Psi\) for the total scalar twist. Take the
strict-at-\(w\) determinant over \(K\) (full elsewhere at allowed finite
places), with old split rational support \(S\) containing 2 and the primes of
\(N_E m\); additionally allow the \(r_j\). Use this both on
\(T_2E\otimes\mathcal O\) and on
\[
                     T_1\oplus T_2=\mathcal O(1)\rho\oplus
                                             \mathcal O\rho^{-1}
\]
Here \(T_2\) standing alone in the comparator denotes its second summand, and
\(\rho\) restricts by norm.  In both coefficient systems we include the scalar
twist \(\Psi\). Write \(D_E,D\) respectively for the determinants. The strict
complex has square amplitude \(1,2\), including on each summand (no residual
global invariants on either side, and usual global and \(w\)-Euler formulas).
The residual determinants agree up to unit by identical diagrams.

Use also the classical weight-two primitive Eisenstein form of level \(m^2\), trivial nebentypus and coefficients
\[
                        a_l/l=\sum_{b e=l}\rho(b)\rho^{-1}(e)/e
\]
(zero on multiples of conductor primes), with no constant at \(\infty\). Use
both fully \(S\)-depleted antiderivatives (opposite tame orientations), and
for it we need just weighted CM disk-center traces on setting \(t=0\); take
order conductor \(s=\prod r_j\) at the stages. Write \(B=b^+ b^-\) for this
\(t=0\) product and \(B_E\) for the curve's paired measure including \(t\).
Then \(B_E(0,\mathbf u)\equiv B\bmod2\) over extended constants as at
\eqref{eq:M1}--\eqref{eq:M2}. Indeed use common split data and level raising
for the quotient operators, the same class-label translates and the canonical
connected choices. The ordinary modular expansion proof there applies equally
to this holomorphic modular form; its fully depleted primitive has Tate
coefficients \(a_l/l\) off the support, integral and congruent as required by
good residual Frobenius and Hecke recursion, and zero elsewhere. The
comparator Euler polynomials in the tame operators and in the log formula
below are
\[
       P_q(Z)=(1-\rho(q)Z)(1-\rho^{-1}(q)Z/q),
                 \qquad P_q=1\ \text{if }q\mid m.
\]

\subsection{Residual regularity and horizontal divisibility}

\begin{lemma}[Addressed conductor variables]\label{cyclic:addressing}
The inert prime sequences and order-conductor variables can be chosen
so that
\[
 \rho(r_j)=1,\qquad a_{r_j}(E)\longrightarrow a_{j,*}\ne0,
 \qquad D\big|_{2=t=0}\ne0.
\]
Moreover, the old split places have nonzero Frobenius exponents, and
each new place has a primitive relative inertia exponent.
\end{lemma}

\begin{proof}
\emph{Addressing the old places.}  Begin with a prime-sequence \(r_0\) giving
nonzero \(u_0\)-Frobenius exponents at the two places of each old \(q\in S\).
Indeed use the power residue symbols at \(r_0\) of 2-power order tending to
infinity dividing \(r_0+1\) (they kill rational residues, and give primitive
relative inertia; the kernel of lowering order and the odd class number give
the conductor quotients). On principal generators of the class-number powers
of one place at each \(q\), nonzero exponents can be forced by taking rational
Frobenius lifts \(\gamma_0\) near \(c\eta\), with \(\eta\in G_K\) fixing
\(\rho,\mu_{2^\infty}\).

Their squares read the power symbols as in the preceding anchors, giving
differences of Kummer translations on those numbers and their conjugates.
These numbers are independent by valuations, and restriction to the specified
kernel retains rational Kummer independence (central cyclotomic scalar); hence
the tuple translation image is open. We may preserve all differences nonzero
and also force the nonzero elliptic limiting trace by varying over the kernel
of these extra radicals (non-CM open image with competing data of bounded
derived length as in the evaluation comparison). Apply Chebotarev with
increasing precisions.

\par\smallskip
\emph{Eliminating residual cohomology.}  After each step in the generic
residual field at \(t=0\), the twisted \(T_1,T_2\) strict \(H^1\)'s have equal
dimension: fixed odd locals are acyclic by the Frobenius exponents or inertia,
moving locals by nonconstant inertia, so the two problems are exact conjugate
duals up to shift and have Euler characteristic zero.

When this dimension is nonzero, arrange nonzero Frobenius-square evaluations
of one class from each at the next prime (start there with \(u_j=0\), old
problem via unramified inflation). The finite lines switch to transverse
singular lines, dropping dimension by at least one on each (old opposite
finite evaluations kill any new transverse contribution by duality).

Here the local limit is split on each summand residually by the tame/Frobenius
calculation: with the new scalar as well it has inertia difference
\((1+u_j)^{b_j}-1,\ b_j\in\mathbb Z_2^\times\), Frobenius difference zero
(\(\gamma_j^2\) trivial on ring fields), residue conjugating power tending to
one. Thus the pure lines at this specialization have the exact cross
orthogonality of the inert switch, and the singular condition including degree
zero is the specialization of the isolated inertia block in degrees \(0,1\).
Off \(u_j=0\) generically it and full give the same cohomology. The dimension
drop therefore persists as an upper bound in the new generic field by the
finite free models.

\par\smallskip
\emph{Realizing the evaluations.}  For these choices use \(\gamma_j\)
approaching \(c\eta_i\) at the stages \(i\), where \(\eta=(\eta_i)\) belongs
to the product kernel over \(K\) of ring-class and cyclotomic data and
\(\rho\). Indeed the classes above inject into unrestricted crossed classes
(notably \(H^0(K_w)=0\) in that residual test) by the evaluation lemma. Their
restrictions are detected on this kernel since its product quotient is abelian
with a nonidentity scalar.

For each line the evaluations on \((c\eta)^2\) are additive and span on the
tested class: conjugate kernel evaluations cannot identically cancel the
original ones as they transform under distinct characters (inverse varying
scalar, already infinite order via \(u_0\)). One can thus achieve both
nonvanishings (union of two proper kernels).

Perturb by products of squares there preserving the test, to force nonzero
elliptic limiting trace: a uniform open determinant-one piece of the curve's
Tate image is available on that kernel as above, also after generating by
squares. Chebotarev with the prescribed evaluations now works. This proves the
desired arrangements by iteration.
\end{proof}

\begin{lemma}[Curve divisibility]\label{cyclic:divisibility}
With these choices, \(D_E\mid B_E\) in the integral power-series ring.
\end{lemma}

\begin{proof}
For the curve, the integral divisibility \(D_E\mid B_E\) now holds by the
high-\(\theta\) horizontal argument \eqref{eq:B1}, \eqref{eq:D2} and residual
regularity. To reiterate the hypotheses here: the added inert locals are
contractible on characteristic-zero tests by nonzero limiting Tate traces on
the rational Frobenius lifts (determinant tending to \(-1\)), and the fixed
odd singular factors and the two Kummer log lines give \eqref{eq:B1} by
\eqref{eq:M2}, the nonzero Frobenius-exponent substitution, and high dyadic
inertia as in the several-conductor-variable argument.

Local fields of the Heegner orbits contain the prescribed finite ramified
character field and growing unramified layers, with bounded ramification at
fixed depth. The switches of \eqref{eq:D2} use additional derivative primes as
in \eqref{eq:K} retaining these base conductors, with Tate absolute
irreducibility over \(K\) by non-CM open image and distinct inverse-twist
determinants on the high-character fibers. Thus the determinant valuation test
and Weierstrass remainder proof there apply without a residual absolute
simplicity hypothesis.
\end{proof}

Take \(t=0\) and a group-power tame line \(1+u_j=(1+v)^{n_j}\) retaining
residual determinant nonvanishing and nonzero exponents on each old Frobenius
and each relative inertia. The integral series tests allow all these
conditions. For all fixed sufficiently high roots \(1+v\) of 2-power order we
will prove for the comparator
\[
                               v_2(B)=v_2(D).                         \tag{Y1}\label{eq:Y1}
\]
It suffices to prove this at actual stages large on the ultrafilter, with the
resulting fixed-order scalar character denoted by \(\chi\): the determinant
valuations stabilize by residual regularity and Weierstrass, and bounded
measure evaluations pass at these fixed scalar precisions as well. The stage
strict determinants are in particular nonzero. The character is unramified at
old support with high order on those Frobenius values, and is nontrivial with
high order on inertia at every added prime; \(\gamma_j^2\) acts trivially on
it.


\subsection{Primitive ray coordinate of the comparator}

Fix such a stage and character. Write \(S'\) for the finite allowed places of \(K\), \(s=\prod_j r_j\), and
\[
             \tau=\rho\chi,\qquad \mathfrak f=(m s),\qquad
             e_S=\prod_{\mathfrak q\mid S,\ \mathfrak q\nmid m}
                               (1-\tau(\mathrm{Fr}_{\mathfrak q})).
\]
Write \(\mathcal O(\tau)\) for the coefficient ring obtained from \(\mathcal
O\) by adjoining the values of \(\tau\).  Here \(\mathfrak f\) is the exact
conductor by split primitivity at \(m\) and nontriviality at the new inert
primes. Let \(z_{\rm ray}\in H^1(K,\operatorname{Frac}(\mathcal
O(\tau))(1)\tau)\), with support \(S'\), be the unnormalized weighted Kummer
trace from the ray field of \(\mathfrak f\) using a Siegel function \(g\) at a
primitive \(\mathcal O_K/\mathfrak f\)-division generator of a maximal-order
CM elliptic curve. We use the standard Siegel function \(g_{a,b}\) (point
\(a\tau_{\rm mod}+b\) in modular-parameter coordinates), leading \(q_{\rm
Tate}^{B_2(a)/2}\) up to phases in its product for \(0<a<1\); \(B_2\) is the
second Bernoulli polynomial. Take a power defining the modular unit without
root ambiguity, then divide its rational Kummer class by that power. We use
the classical Siegel-function transformation laws and the main theorem of CM
(see \cite[Section~2.2]{Oukhaba03}): on changing bases or acting on the
elliptic curves with the point, these functions transform with the torsion
labels up to roots of unity (on determinant actions the Tate phases transform
as the roots of unity). Thus the powers can be taken in the ray field, with
ray classes acting by the ideal isogenies, and replacing the generator/base
label only translates the sum or changes killed roots.

The unit powers have integrality and invertibility off level divisors at
integral CM \(j\)'s. Equivalently this follows by taking symmetric polynomials
of all modular conjugates (also for inverses), with no interior poles, and
using the Siegel products at cusps to check integrality of the polynomials in
\(j\) away from the level.

\begin{proposition}[Primitive ray determinant]\label{cyclic:ray-basis}
The class
\[
 z_S=e_Sz_{\rm ray}
\]
is a basis of the inverse determinant of the unrestricted global
complex on \(\mathcal O(\tau)(1)\tau\), with support \(S'\), over the
dyadic character integers.  Rationally this determinant is its single
\(H^1\)-line.
\end{proposition}

\begin{proof}
\emph{Regular representations and the class-number formula.}  First do the
determinant and regulator comparisons for regular representations in the
cyclic tower over \(K\) cut out by \(\tau\). At an intermediate field \(F\)
these use by Shapiro the \(2\)-adic \(\mathbb Z_2(1)\) complex over \(F\) for
\(S'\) inverted. Its \(H^1\) is compact \(S'\)-units, \(H^2\) is given by
\(S'\)-ideal classes (2-part) and the sum-zero lattice
\(Y_{S'}^0\otimes\mathbb Z_2\) of invariants, and other cohomology vanishes
(Kummer, Brauer reciprocity). Here \(Y_V\) denotes the label lattice on places
of \(F\) in a set \(V\). Thus relative to free unit and invariant bases the
generator of the inverse determinant has coefficient of valuation
\(v_2(h_{S'}/w_F)\) (class number over torsion). Compare now by the ordinary
real regulator
\[
        \operatorname{units}_{S',F}\otimes\mathbb R
                    \ \simeq\ (Y_\infty\oplus Y_{S'})^0_{\mathbb R},
\]
taking \(-2\log|\cdot|\) at complex labels (one embedding per label extending
the chosen embedding of \(K\)) and valuations times log norms at finite ones.
Cancel the degree-two invariant coordinates by the exact sequence of
label lattices
\[
 0\longrightarrow Y_{S'}^0\longrightarrow
 (Y_\infty\oplus Y_{S'})^0\longrightarrow Y_\infty
 \longrightarrow0.
\]
The last map is projection to the infinite labels. It is surjective
because \(S'\ne\varnothing\): any total degree at infinity can be
balanced at one finite label. No choice of a splitting enters the
determinant comparison. Thus an
inverse-determinant tensor \(r\) on rational cohomology bases and a basis
tensor \(b_\infty\) on \(Y_\infty\) have signed comparison factor
\(j(r,b_\infty)\) (coefficient of the image of \(r\)). By the classical
\(S'\)-truncated Dedekind class number formula at zero, leading zeta value
divided by this factor is rational, of valuation
\[
                     v_2(b_\infty\text{ relative to lattice basis})
                                      -d(C_F,r).
\]
Indeed on ordinary free unit, invariant and label lattice bases this is just
\(h_{S'}/w_F\) up to sign after cancelling the regulator.

Under intermediate field restriction all these comparisons commute: normalized
logs and Brauer invariants extend with local degree multipliers, at infinity
just unweighted injection by subgroup sums on labels.

\par\smallskip
\emph{The integral primitive quotient.}  Quotient to the primitive powers of
\(\tau\), whose order is \(3\cdot2^a\) say. First remove the part trivial on
the cubic subgroup by subgroup-sum injection; in the remaining module remove
the \(2^{a-1}\)-part by the analogous injection (polynomial \(1+X^{2^{a-1}}\)
in the 2-power factor). Thus the quotient lattice is the cyclotomic integers
with group map
\(g\mapsto\tau(g)\).  More explicitly, put
\(G=C_3\times C_{2^a}\), let \(N_3\) be the subgroup sum for
\(C_3\), and let \(X\) generate the two-power factor.  The two
coefficient sequences are
\[
\begin{gathered}
0\longrightarrow\mathbb Z_2[C_{2^a}]
 \xrightarrow{\ N_3\ }\mathbb Z_2[G]
 \longrightarrow\mathbb Z_2[\mu_3][C_{2^a}]
 \longrightarrow0,\\
0\longrightarrow\mathbb Z_2[\mu_3][C_{2^{a-1}}]
 \xrightarrow{\ 1+X^{2^{a-1}}\ }
 \mathbb Z_2[\mu_3][C_{2^a}]
 \longrightarrow\mathbb Z_2[\mu_{3\cdot2^a}]
 \longrightarrow0.
\end{gathered}
\]
The high-order character ensures \(a\ge1\). 

These sequences on regular modules and their successive quotients give
integral determinant triangles, with the subgroup-sum maps using restriction
under Shapiro. Rationally the same sequences identify the quotient regulator
comparisons by the preceding compatibility, including on infinity lattices by
the subgroup sums. Hence by multiplying/dividing the preceding equalities we
have the same determinant-ratio valuation rule for this primitive part, using
as zeta leading value the product of \(L'_{S'}(\tau^i,0)\) on primitive
powers. Indeed these characters have a simple zero individually (Hecke
functional equations and central-side-at-one nonvanishing by Dedekind
factorization for the finite characters), with no extra zero from truncation.

Rationally in cohomology only the character-field unit line in degree one
remains, since there are no corresponding \(S'\)-invariant coordinates
(\(\tau\) nontrivial at every allowed place). On infinity labels, in any
primitive embedding the map of a pushed-forward ordinary unit to the regulator
coordinate uses
\[
                                -2\sum_g\tau(g)\log|g(u)|
\]
with the corresponding embedded character. Indeed index the regular labels by
embedding composed with \(g\); Shapiro inverse corestricts with identity-label
basis, and group-ring multiplication acts inversely on the unit. Norming first
from the ray field therefore uses exactly such full sums over the ray labels
on \(z_{\rm ray}\).

\par\smallskip
\emph{The ray coordinate and the limit formula.}  Each primitive embedding on
\(z_S\) now gives precisely a primitive \(L'_{S'}(0)\)-coordinate, up to
roots/signs and consistent permutation of embeddings. Indeed each ray partial
zeta for an integral prime-to-\(\mathfrak f\) representative \(\mathfrak a\)
is a norm scale times the shifted sum of \(|x|^{-2s_1}\) over \(1+(m
s)\mathfrak a^{-1}\); there is no nontrivial global unit congruent to one. Its
value at \(s_1=0\) vanishes and derivative is \(-2\log|g|\) at the shift
generator on the corresponding maximal-order lattice. This is the shifted
Epstein--Kronecker limit formula.  To check its
normalization against the ray invariant, write \(\mathfrak f\cap\mathbb Z=f_{\mathbb Z}\mathbb Z\) and take its defining power
\(N=12f_{\mathbb Z}\).  In the notation of
\cite[Section~2.2 and Equation~(2.4)]{Oukhaba03}, the invariant is
\(g^N\) up to a phase, and the congruence-unit number
\(w_{\mathfrak f}\) is one here.  The factor
\(-N^{-1}\log|g^N|^2\) is therefore exactly \(-2\log|g|\).
The rational Kummer class was divided by this same defining power. For example
after scaling to lattice \((\tau_{\rm mod},1)\) move the shift to \(a\tau_{\rm
mod}+b,\ 0<a<1\); Poisson summation in the second coordinate gives in the
derivative the constant mode \(2\pi\operatorname{Im}(\tau_{\rm mod})B_2(a)\)
and other modes
\[
  \sum_{n\in\mathbb Z,\,l\ge1}
    \frac2l e^{-2\pi l|n+a|\operatorname{Im}(\tau_{\rm mod})}
                  \cos(2\pi l(b+(n+a)\operatorname{Re}(\tau_{\rm mod}))),
\]
exactly that product-log. Varying \(\mathfrak a\) uses the ray conjugates by
CM reciprocity, and \(e_S\) supplies the extra omitted Euler terms.
Ideal-action inversion conventions at most permute primitive powers here and
replace Euler omissions by their inverse-character versions, of ratio a root
of unity.

Use the tensor formed from \(z_S\) times a \(\mathbb Z\)-basis of the
cyclotomic-integer lattice as \(r\), and the same lattice basis on infinity as
\(b_\infty\). These are rational comparisons (units from powers before
dividing, not a rational structure inferred from real logs). Their real
determinant comparison therefore has exactly the absolute value of the
primitive leading-value product. This gives \(d(C_{\rm prim},r)=0\) over
\(\mathbb Z_2\). There is only one prime above 2 in the coefficient field
\(\mathbb Q(\mu_{3\cdot2^a})\). Let \(F_2\) be its dyadic
completion and write \(z_S=a e\) for a generator \(e\) of the
character determinant lattice. On restriction of scalars the index
is the norm of \(a\), and
\[
 v_2\bigl(N_{F_2/\mathbb Q_2}(a)\bigr)
                  =[F_2:\mathbb Q_2]v_2(a).
\]
The norm-index valuation just computed is zero, so \(v_2(a)=0\).
This proves the individual determinant-basis assertion. Further
coefficient enlargements preserve it.
\end{proof}

\subsection{The local logarithm lattice and the strict determinant}

Write \(\ell_w\) for local untwisted log after restriction to the local
unramified character field \(P_w/\mathbb Q_2\) of \(\tau\) and evaluation in
the chosen embedding. Strict nonvanishing already arranged thus gives \(\ell_w
z_S\ne0\). Put
\[
 a_w=v_2(1-\tau(\mathrm{Fr}_w))=v_2(1-\tau(\mathrm{Fr}_{\bar w}))\le1
\]
(using the high Frobenius-power test).

\begin{lemma}[Unramified logarithm lattice at two]\label{cyclic:local-log}
The evaluated free logarithm lattice of
\(H^1(K_w,\mathcal O(\tau)(1)\tau)\) has valuation \(1+a_w\).
Its torsion and the local \(H^2\) each have normalized length \(a_w\).
Consequently, the strict determinant on this summand has valuation
\[
 v_2(\ell_wz_S)-1-a_w.
\]
\end{lemma}

\begin{proof}
Put \(A=\mathcal O(\tau)\) and \(P=P_w\).  Restriction identifies the
local \(H^1\) integrally with the twisted invariants of
\(P^{\times,\wedge}_2\otimes A\): the integral \(H^0\) of the Tate
twist upstairs is zero.  The valuation subgroup contributes no
invariants.  The sign torsion contributes normalized length \(a_w\),
and local duality gives the same length for \(H^2\).

Set \(\Gamma=\operatorname{Gal}(P/\mathbb Q_2)\),
\(U^1=1+2\mathcal O_P\), and \(U^2=1+4\mathcal O_P\).
On tensor products with \(A\), let \(\Gamma\) act diagonally,
using the character twist on \(A\). The logarithm identifies
\(U^2\) with \(4\mathcal O_P\). An integral normal basis makes
this an induced \(\Gamma\)-lattice, also after the character twist;
in particular \(H^1(\Gamma,U^2\otimes A)=0\). Thus
\[
\begin{split}
0\longrightarrow (U^2\otimes A)^\Gamma
 &\longrightarrow (U^1\otimes A)^\Gamma\\
 &\longrightarrow ((\mathcal O_P/2)\otimes A)^\Gamma
 \longrightarrow0
\end{split}
\]
is exact. The last term is isomorphic as an \(A\)-module to
\(A/2A\), so has normalized length one. Evaluation of the smaller
logarithm lattice has valuation exactly two: the character resolvent
of an integral normal generator is a unit, since the normal-basis
circulant matrix is invertible over the integers.

The sign torsion, of length \(a_w\), is killed by the logarithm and
has trivial intersection with \(U^2\).  The free logarithm lattice is
therefore enlarged from the smaller lattice by length \(1-a_w\).
Its valuation is \(2-(1-a_w)=1+a_w\).  Finally, in the global--local
determinant triangle the local torsion terms have equal lengths and
cancel.  Proposition~\ref{cyclic:ray-basis} then gives the stated
strict determinant valuation.
\end{proof}

\begin{proposition}[Comparator determinant]\label{cyclic:determinant}
At every stage and high-order character selected above,
\[
           v_2(D)=2v_2(\ell_w z_{\rm ray})+
                  \sum_{q\in S,\,v_0\mid q}
                             v_2(P_q(\chi(\mathrm{Fr}_{v_0}))).
                  \tag{Y2}\label{eq:Y2}
\]
The two terms above the rational prime two sum to \(2a_w-2\).
\end{proposition}

\begin{proof}
The conjugate dual strict problem has the same valuation by duality. Passing
back from its zero to full odd allowed places on \(T_2\chi\) adds at each
unramified old odd place \(v_0\mid q\) the length difference
\[
                       v_2(1-q/\beta)-v_2(1-\beta),\qquad
                  \beta=(\rho^{-1}\chi)(\mathrm{Fr}_{v_0}),
\]
by \(H^2,H^1\) torsion respectively. 

There is integral acyclicity at primes of \(m\) by cubic inertia. At each new
inert place the difference is zero: \(H^1\)-torsion there uses discrete
invariants on the coefficient modulo its lattice, with Frobenius lift
\(\gamma_j^2\) acting trivially, hence has length the valuation of inertia
scalar minus one; the dual \(H^2\)-length is the same (\(r_j^2\) sufficiently
close to 1 here). These complexes are all rationally acyclic by the character
tests. Since \(v_2(1-\tau(\mathrm{Fr}_{v_0}))=v_2(1-\beta)\), adding
these terms to Lemma~\ref{cyclic:local-log} and substituting
\(z_S=e_Sz_{\rm ray}\) gives \eqref{eq:Y2}.  At two, each polynomial
has valuation \(a_w-1\): its first factor has valuation \(a_w\),
and its second factor has valuation \(-1\).  The two places
therefore contribute \(2a_w-2\), as asserted.
\end{proof}


\subsection{Comparator disk centers}

The primitive log to be used for the Eisenstein form in the disk formula is
\[
 G=-\mathfrak g(\rho)^{-1}
        \sum_{\alpha,\beta\in(\mathbb Z/m)^\times\ /\ {\rm simul.}\ \pm}
           \rho(\alpha\beta)\log g_{\alpha/m,\beta/m}(m\tau_{\rm mod}),
 \qquad \mathfrak g(\rho)=\sum_b\rho(b)e^{2\pi i b/m}.
\]
\begin{lemma}[Eisenstein primitive]\label{cyclic:eisenstein-primitive}
The function \(G\) has Tate coefficients \(a_l/l\), and its
three degree-two quotient logs sum to \(a_2G\).  Hence
\(P_2(V_2)G\) is the two-depleted antiderivative on the ordinary
CM disks.
\end{lemma}

\begin{proof}
Transfer constants accordingly to the embedding used; the Gauss factor is a
2-adic unit. Evaluate thus on the middle curve of the cyclic \(m^2\)-chain,
the two axes being generators of the dual-previous and next subgroups with a
fixed Weil pairing. Inverse scalings make no difference. The log Tate
expansion at the cyclic \(\mu\)-level cusp has exactly coefficients \(a_l/l\):
the two opposite progressions in each Siegel product with the simultaneous
sign quotient retain one copy of each progression, then sum by Gauss over
\(\beta\). Leading \(q_{\rm Tate}\)-powers cancel since
\(\sum_\beta\rho(\beta)=0\); root phases have log zero.

In particular the differential agrees with the prescribed modular-form
differential. Also the sum of logs at the three degree-two quotients is \(a_2
G\), with no constant. Indeed at the Tate cusp use the two square-root
expansions and square expansion giving this by Hecke recursion. The
differential identity extends algebraically, and this constant-log comparison
is of modular-unit combinations: zero differential gives a zero divisor
combination, hence a scalar combination of integral divisor relations,
reducing to logs of constants checked at the same cusp. One can pull to a fine
cover with the relevant branches for this test, taking powers to remove root
ambiguities. The logs are analytic on the ordinary interior disks. Thus the
two-source trace argument of \eqref{eq:M1} applies, and \(P_2(V_2)G\) is
precisely the 2-depleted antiderivative there.
\end{proof}

At the stage character this gives the two primitive weighted order-conductor
sums of \(G\), multiplied by the Euler polynomials in the inverse translations
on the respective orientations. The polynomial-factor valuation sum is exactly
the sum in \eqref{eq:Y2} (opposite choices at odd primes; at 2 repeating the
connected choice doesn't change the valuation).

\begin{lemma}[Order labels and ray labels]\label{cyclic:ray-labels}
In either orientation, the primitive weighted order-conductor sum of
\(G\) has valuation \(v_2(\ell_wz_{\rm ray})\).
\end{lemma}

\begin{proof}
The middle-curve passage translates the order-\(s\) labels by a fixed split
ideal. The \((\alpha,\beta)\) yield every primitive \(m\)-division generator
over the order modulo sign, since the axes are split. Under actual CM Galois
action over \(K\) pairing-normalized frames on the axes transform by scalings
of product the mod-\(m\) cyclotomic value; thus the weights
\(\rho(\alpha\beta)\) transform by exactly the cubic norm character.

Parametrize the middle curves of order conductor \(s\) by maximal-order curves
quotiented by one line at each inert \(r_j\), transporting the split groups.
This gives all the labels exactly by CM and the order class group (relative
transitivity on all the lines, units just signs). Pairing-frame transport from
below has here just a fixed extra degree scaling from \(s\). Through these
isogenies the Siegel log sums over all preimages of the torsion point, without
averaging (for a cyclic quotient multiply the products for the Tate
multiplicative kernel, and use basis change, up to killed phases).

Terms omitting any \(r_j\)-component trace to zero by nontriviality of
\(\chi\) on its line labels. In the surviving terms, each nonzero
\(r_j\)-division point lies on exactly one of the \(r_j+1\) lines.  Each line
supplies \(r_j-1\) nonzero points, giving \((r_j+1)(r_j-1)=r_j^2-1\) labels.
These are exactly the primitive \(\mathcal O_K/(r_j)\)-coordinates by
inertness.

The label identification can be made explicitly on each maximal-order
CM curve. Choose representatives for the \(m\)-division generators
modulo simultaneous sign. The Chinese remainder decomposition combines
such a generator with one nonzero \(r_j\)-division point for every
\(j\) to give a primitive \(\mathcal O_K/(ms)\)-coordinate modulo
units. Conversely, each nonzero \(r_j\)-coordinate determines its
unique \(\mathbb F_{r_j}\)-line, recovering the quotient-curve
label. Changing the chosen \(m\)-representative negates all the
\(r_j\)-coordinates simultaneously. Thus there is exactly one
global sign identification, and these are all the ray labels above
the Hilbert classes.

Their ring-order weights use the natural line-quotient curves, and
with the frame weights thus
give exactly the full \(\rho\chi\) on the labels by CM action, up to
translation (ideals labelled here by actual Artin action). This works on
either orientation. Hence we have up to units precisely the ray log on
\(z_{\rm ray}\), which likewise uses the full sum after restricting weighted
corestriction.

Even primitive conjugates under convention changes have the same valuation:
coefficient automorphisms at the unique dyadic prime extend locally and act on
the underlying Siegel logs by a ray translation (\(K_w=\mathbb Q_2\)).
\end{proof}

Combining Proposition~\ref{cyclic:determinant} with
Lemmas~\ref{cyclic:eisenstein-primitive} and~\ref{cyclic:ray-labels} proves
\eqref{eq:Y1}.

\subsection{The cyclic-cubic anchor}

\begin{proposition}\label{cyclic:anchor}
For the original curve \(E_0\), the discriminants \(h\) and \(-d\)
of Lemma~\ref{cyclic:seed} satisfy
\[
 \an(E_0^h)+\an(E_0^{-dh})=1,\qquad
 X(E_0^h)+X(E_0^{-dh})=0.
\]
Thus Proposition~\ref{prop:anchor} holds for cyclic-cubic residual
image.
\end{proposition}

\begin{proof}
\emph{The integral unit.}  By \eqref{eq:Y1}, the restrictions of
\(B,D\) to the chosen tame line have the same finite residual
Weierstrass order: evaluate at sufficiently high two-power roots and
apply the valuation test.  The residual congruences
\(B_E(0,\mathbf u)\equiv B\pmod2\) and \(D_E\equiv D\pmod2\),
with determinant bases chosen up to units, give the same equality of
orders for the curve.  Lemma~\ref{cyclic:divisibility} makes
\(B_E/D_E\) integral.  Its residual restriction has order zero, so
its constant term is a unit and therefore \(B_E/D_E\) is a unit in
the full local power-series ring.

\emph{The rank-one center.}  Apply the rank-detection argument from the proof of
Proposition~\ref{split-local:center-rank} at \(t=0\), on the transverse tame line through the
characteristic-zero origin.  Lemma~\ref{cyclic:seed} gives
\(s_2(E/K)=1\).  The dyadic comparison is the unramified
Kummer/logarithm comparison, and all other allowed finite terms are
acyclic there; for the moving inert terms this uses
\(a_{j,*}\ne0\).  Formula~\eqref{eq:M2} identifies the center log
square. Together with the unit quotient \(B_E/D_E\), it gives zero
valuation of the finite class-functional tensor.  Both special
cohomology lines have rank one, so the Heegner class specializes
nontrivially.  The exact Gross--Zagier formula \eqref{eq:GZ-E} now
gives a simple zero of \(L(E/K,s)\).  In particular each elliptic
factor has analytic rank at most one, and its discrepancy \(X\) is
defined.

\emph{The exact central value.}  Evaluate the same unit equality at
the all-zero origin.  At every added inert prime, the singular
unramified-inflation determinant contributes valuation
\(2v_2(a_{j,*})\); the two oriented point traces contribute the
same squared multiplier \(a_{j,*}^2\).  They cancel integrally, as
in the preceding anchors.  Formula~\eqref{eq:E}, now with precisely
the support \(S\), and the arithmetic-volume identity
\eqref{eq:G} yield
\[
 X(E)+X(E^{-d})=0.
\]
The real period and Gross--Zagier height are the exact normalizations
of \eqref{eq:GZ}.  Returning to \(E=E_0^h\) gives both asserted
equalities.  Finally, \(h\) is an odd fundamental discriminant prime
to \(2N_{E_0}\), and \(-d\) is prime to \(2hN_{E_0}\), whose prime factors
all split in \(K\).  These are precisely the discriminant
conditions in Proposition~\ref{prop:anchor}.
\end{proof}
